\section{模型与算法介绍}

\subsection{ResNeXt 与 DenseNet}

ResNeXt 是在残差网络基础上发展起来的一类模型，其核心思想是将多条具有相同拓扑结构的变换分支进行聚合，并通过 cardinality 提升表示能力。与单纯增加网络深度或宽度相比，这种做法在保持结构规则性的同时，能够更高效地提升特征表达能力。本文选用的 ResNeXt50\_32x4d 是一个较常见的标准版本，适合与其他经典 CNN 进行对比。

DenseNet 的特点是密集连接，即每一层都会接收前面所有层的特征图作为输入，并将自己的输出继续传递给后续层。这样的设计有利于特征复用，也能够在一定程度上缓解梯度传播过程中的信息损失。本文选用 DenseNet121 作为第二个对照模型，希望观察其在小样本图像分类中的稳定性表现。

\subsection{从头训练与微调}

从头训练是指模型参数全部随机初始化，然后仅依靠当前任务数据进行优化。这种方式的优点是训练过程完全针对目标任务展开，不依赖外部数据；但缺点也比较明显，当数据量较小时，模型往往更依赖超参数设置，训练初期容易震荡，收敛速度通常较慢。

微调则是将模型初始化为在 ImageNet 上训练得到的公开预训练权重，再针对当前任务继续训练。考虑到预训练模型已经学到了较通用的边缘、纹理和中层语义特征，微调通常会在小样本任务中表现出更快的收敛速度和更高的起始性能。为了让对比更公平，本文对两种模型都采用统一的两阶段微调方案：先冻结主干网络，仅训练最后的分类层；随后解冻全部参数，以较小学习率进行全模型优化。

\subsection{训练设置}

四组实验均使用交叉熵损失函数和 AdamW 优化器，输入图像统一缩放到 $224 \times 224$。从头训练的初始学习率设置为 $3 \times 10^{-4}$，训练 30 个 epoch；微调实验先以 $1 \times 10^{-3}$ 训练分类头 5 个 epoch，再以 $1 \times 10^{-4}$ 解冻全模型训练 15 个 epoch。评价指标使用测试集 Accuracy 和 Macro-F1，并记录训练总时长、最优轮次以及训练过程中的曲线变化。
